\section*{Chapitre \ref{ch:b3:solid-state} --- Chimie du solide~: bandes, défauts, semi-conducteurs}

\begin{solution}{exo:b3:solid-state:1}
$(E - \alpha)/|\beta| = -2\cos(k\pi/7)$~: $-1.802$, $-1.247$, $-0.445$, $+0.445$, $+1.247$, $+1.802$. Ils couvrent
$3.604|\beta|$, déjà 90\,\% de la largeur $4|\beta|$ de la bande de la chaîne infinie.
\end{solution}

\begin{solution}{exo:b3:solid-state:2}
$2N$ ($N$ orbitales, deux électrons chacune). Le sodium, avec un électron 3s par atome, remplit à moitié
sa bande s~: c'est un métal. Le magnésium en a deux et remplirait
exactement la bande s, mais les bandes 3s et 3p se chevauchent, si bien que
les niveaux occupés les plus hauts appartiennent à une bande combinée
partiellement remplie~: c'est aussi un métal.
\end{solution}

\begin{solution}{exo:b3:solid-state:3}
$\lambda \approx \qty{1239.8}{eV.nm}/E_g$~: phosphure de gallium \qty{549}{nm} (vert)~; nitrure de gallium \qty{365}{nm}
(proche ultraviolet). Les diodes bleues utilisent du nitrure de gallium
allié à l'indium, dont la \omterm{def:b3:solid-state:band}{bande interdite} est plus étroite.
\end{solution}

\begin{solution}{exo:b3:solid-state:4}
$\ce{CaCl2} \xrightarrow{\ce{NaCl}} \mathrm{Ca_{Na}^{\bullet} + V_{Na}' + 2\,Cl_{Cl}^{\times}}$~: deux sites cationiques et deux sites anioniques, comme dans NaCl~; un Ca et
deux Cl~; charge $+1 - 1 = 0$. $\ce{Y2O3} \xrightarrow{\ce{ZrO2}} \mathrm{2\,Y_{Zr}' + 3\,O_O^{\times} + V_O^{\bullet\bullet}}$~: deux sites cationiques et quatre sites d'oxygène~; deux Y
et trois O~; charge $-2 + 2 = 0$.
\end{solution}

\begin{solution}{exo:b3:solid-state:5}
$T^{3/2} = 5196$ à \qty{300}{K}, $kT = \qty{0.02585}{eV}$. Silicium~: $\sqrt{N_cN_v} = \qty{2.42e19}{cm^{-3}}$, $\eu^{-1.125/0.05170} = \num{3.55e-10}$,
$n_i = \qty{8.6e9}{cm^{-3}}$. Germanium~: $\sqrt{N_cN_v} = \qty{7.16e18}{cm^{-3}}$, $\eu^{-0.661/0.05170} = \num{2.80e-6}$, $n_i = \qty{2.0e13}{cm^{-3}}$. Rapport
$\num{4.3e-4}$~: la \omterm{def:b3:solid-state:band}{bande interdite} plus étroite du germanium lui donne plus
de deux mille fois plus de porteurs.
\end{solution}

\begin{solution}{exo:b3:solid-state:6}
$n = N_D = \qty{1.0e16}{cm^{-3}}$~; $p = n_i^2/n = \num{1.0e20}/\num{1.0e16} = \qty{1.0e4}{cm^{-3}}$.
\end{solution}

\begin{solution}{exo:b3:solid-state:7}
$E_F - E_i = kT\ln(n/n_i) = \qty{0.02585}{eV} \times \ln(10^6) = \qty{0.36}{eV}$.
\end{solution}

\begin{solution}{exo:b3:solid-state:8}
$kT = \qty{0.06894}{eV}$ à \qty{800}{K}~; $n/N = \eu^{-2.3/0.1379} = \num{5.7e-8}$~; dans une mole, $\num{6.02e23} \times \num{5.7e-8} = \num{3.4e16}$
paires.
\end{solution}

\begin{solution}{exo:b3:solid-state:9}
La masse d'une maille vaut $\rho a^3 = \qty{5.70}{g.cm^{-3}} \times (\qty{4.300e-8}{cm})^3 = \qty{4.532e-22}{g}$, soit
$\qty{4.532e-22}{g} \times N_A = \qty{272.9}{g}$ par mole de mailles, ou \qty{68.23}{g} par oxygène. Alors
$(1 - x) \times 55.8 + 16.0 = 68.23$ donne $1 - x = 0.936$~: $x = 0.064$, $\mathrm{Fe_{0.936}O}$.
\end{solution}

\begin{solution}{exo:b3:solid-state:10}
$\ln(\sigma T)$~: $-3.51$, $-1.30$, $0.36$ pour $1/T = 3.333$, $2.857$, $\qty{2.500e-3}{K^{-1}}$~; pente entre les points extrêmes
$(0.36 + 3.51)/(-\qty{8.33e-4}{K^{-1}}) = -\qty{4.65e3}{K}$ (le point du milieu est sur la droite), donc
$E_a = \qty{4.65e3}{K} \times \qty{8.617e-5}{eV.K^{-1}} = \qty{0.40}{eV}$.
\end{solution}

\begin{solution}{exo:b3:solid-state:11}
$\tfrac12\ce{O2} \rightleftharpoons \mathrm{O_O^{\times} + V_M' + h^{\bullet}}$, $K = [\mathrm{V_M'}][\mathrm h^\bullet]/p(\ce{O2})^{1/2}$. Bilan des charges $[\mathrm h^\bullet] = [\mathrm{V_M'}]$, donc
$[\mathrm h^\bullet]^2 = K\,p(\ce{O2})^{1/2}$ et $[\mathrm h^\bullet] \propto p(\ce{O2})^{1/4}$~; la conductivité suit les \omterm{def:b3:solid-state:carriers}{trous}.
\end{solution}

\begin{solution}{exo:b3:solid-state:12}
$n/N = \sqrt{N_i/N}\,\eu^{-\Delta H_F/2kT} = \sqrt 2\,\eu^{-1.4/(2 \times 0.05170)} = 1.414 \times \num{1.32e-6} = \num{1.9e-6}$.
\end{solution}

\begin{solution}{pb:b3:solid-state:1}
\textbf{1.} $\ce{Y2O3} \xrightarrow{\ce{ZrO2}} \mathrm{2\,Y_{Zr}' + 3\,O_O^{\times} + V_O^{\bullet\bullet}}$.

\textbf{2.} Sites~: deux sites cationiques et quatre sites d'oxygène (trois occupés, un vide),
le rapport 1\,:\,2 de $\ce{ZrO2}$. Matière~: 2 Y et 3 O de chaque côté. Charge~: $2 \times (-1) + 2 = 0$.

\textbf{3.} Une lacune pour deux ions yttrium, soit une demi-lacune par ion yttrium.

\textbf{4.} Pour $(\ce{ZrO2})_{0.92}(\ce{Y2O3})_{0.08}$, il y a 0.92 Zr et 0.16 Y, soit 1.08 cation~: $x = 0.16/1.08 = 0.148$,
$\mathrm{Zr_{0.852}Y_{0.148}O_{1.926}}$.

\textbf{5.} $(x/2)/2 = 0.0370$~: 3.7\,\% des sites d'oxygène sont vides.

\textbf{6.} $8 \times 0.0370 = 0.296$ lacune par maille.

\textbf{7.} $a^3 = (\qty{5.14e-8}{cm})^3 = \qty{1.358e-22}{cm^3}$~: $0.296/\qty{1.358e-22}{cm^3} = \qty{2.2e21}{cm^{-3}}$.

\textbf{8.} Chaque ion oxyde doit franchir une barrière de \qty{1.0}{eV}
pour sauter dans une lacune voisine~; la fraction des tentatives qui réussissent, $\eu^{-E_a/kT}$, croît très vite avec $T$.

\textbf{9.} À \qty{1073.15}{K}, $kT = \qty{0.09248}{eV}$~: $A = \sigma T\eu^{E_a/kT} = 0.030 \times 1073.15 \times \eu^{10.81} = \qty{1.6e6}{S.K.cm^{-1}}$.

\textbf{10.} À \qty{973.15}{K}~: $\sigma = (A/T)\eu^{-11.92} = \qty{0.011}{S.cm^{-1}}$.

\textbf{11.} À \qty{573.15}{K}~: $\sigma = (A/T)\eu^{-20.25} = \qty{4.5e-6}{S.cm^{-1}}$.

\textbf{12.} $R = L/(\sigma S) = \qty{0.10}{cm}/(\sigma \times \qty{0.20}{cm^2})$~: \qty{46}{\ohm} à \qty{700}{\celsius}, $\qty{1.1e5}{\ohm}$ à \qty{300}{\celsius}.

\textbf{13.} $R < \qty{1}{k\ohm}$ exige $\sigma > \qty{5.0e-4}{S.cm^{-1}}$~; la résolution de $(A/T)\eu^{-E_a/kT} = \num{5.0e-4}$ (par tâtonnement)
donne $T \approx \qty{760}{K}$, environ \qty{490}{\celsius}.

\textbf{14.} Les gaz d'échappement sont froids après un démarrage et à
faible charge~; l'élément chauffant porte l'électrolyte à sa température de
fonctionnement en quelques secondes, si bien que la régulation fonctionne
dès le démarrage, quand la plupart des polluants sont émis.

\textbf{15.} $\ce{O2 + 4e- -> 2O^2-}$~; en \omterm{def:b3:solid-state:kroger-vink}{notation de Kröger--Vink}, $\ce{O2} + \mathrm{2\,V_O^{\bullet\bullet} + 4\,e'} \longrightarrow \mathrm{2\,O_O^{\times}}$.

\textbf{16.} Le dioxygène est réduit du côté de l'air et les ions oxyde sont
réoxydés en dioxygène du côté des gaz d'échappement~: la cellule transfère $\ce{O2}$ de l'air vers les gaz d'échappement, avec
$\Delta_rG = RT\ln(p_{\mathrm{exh}}/p_{\mathrm{air}})$ par mole de $\ce{O2}$ et quatre électrons~: $E = -\Delta_rG/4F = (RT/4F)\ln(p_{\mathrm{air}}/p_{\mathrm{exh}})$.

\textbf{17.} $RT/4F = 8.3145 \times 973.15/(4 \times 96485) = \qty{0.02096}{V}$.

\textbf{18.} Zéro~: les deux pressions sont égales.

\textbf{19.} $(RT/4F)\ln 10 = \qty{48.3}{mV}$ par décade.

\textbf{20.} $\qty{0.02096}{V} \times \ln(0.2095/0.010) = \qty{0.02096}{V} \times 3.04 = \qty{0.064}{V}$.

\textbf{21.} $\qty{0.02096}{V} \times \ln(0.2095/\num{1.0e-20}) = \qty{0.02096}{V} \times 44.49 = \qty{0.93}{V}$.

\textbf{22.} Les gaz d'échappement riches contiennent du monoxyde de
carbone, du dihydrogène et des hydrocarbures imbrûlés~; sur l'électrode de
platine, le peu de dioxygène restant réagit avec eux, et $p(\ce{O2})$ est fixée par les équilibres $\ce{2CO + O2 <=> 2CO2}$ et $\ce{2H2 + O2 <=> 2H2O}$, qui laissent une valeur
infime à \qty{700}{\celsius}.

\textbf{23.} $p = 0.2095\,\eu^{-0.45/0.02096} = \qty{1.0e-10}{bar}$.

\textbf{24.} Autour du rapport stœchiométrique, les gaz d'échappement
passent d'un léger excès de dioxygène à un léger excès de CO et de
dihydrogène~: $p(\ce{O2})$ chute de quelque dix-huit puissances de dix pour une faible variation de
carburant, et la tension saute de moins de \qty{0.1}{V} à environ
\qty{0.9}{V}. Le calculateur du moteur ajoute du carburant quand la tension
est basse et en retire quand elle est haute, ce qui maintient le mélange au
point stœchiométrique où le catalyseur trois voies convertit à la fois CO,
les hydrocarbures et les oxydes d'azote.

\textbf{25.} Dans des gaz d'échappement riches à \qty{700}{\celsius}, la sonde délivre environ \qty{0.93}{V}.
\end{solution}
